\documentclass[a4paper, 12pt]{article}

%%%%%%%%%%%%%%%%%%%%%%%%%%%%%%%%%%%%%%%%%%%%%%%%%%%%%%%%%%%%%%%%%%%%%%%%%%%%%%%%

% Math, plots and sketches
\usepackage{amsmath, mathtools, amsfonts}
\usepackage{tkz-euclide}
\usepackage{pgfplots}
\pgfplotsset{compat=1.18}

% Images
\usepackage{graphicx}

% Allow line-breaking URLs/paths in text
\usepackage{url}
\def\UrlBreaks{\do\/\do-\do\_\do\#\do\0\do\1\do\2\do\3\do\4\do\5\do\6\do\7\do\8\do\9}

% Formatting
\usepackage[left=1in, right=1in, top=1.2in, bottom=1.2in]{geometry}
\usepackage{dirtytalk} % Proper quotation marks
\usepackage{float} % Make it possible for tables (and images) to be placed in an explicit location (using [H])


% Colors for highlighting and coloring tables
\usepackage{color, soul}
\usepackage[table]{xcolor}

% Multicolumn formatting
% In case that's the format we want
%\usepackage{multicols}

% Custom labeled lists
\usepackage{enumerate}

% Clickable references
\usepackage{hyperref}
% Source - https://tex.stackexchange.com/a/51349
% Posted by Canageek, modified by community. See post 'Timeline' for change history
% Retrieved 2026-03-23, License - CC BY-SA 4.0
\hypersetup{
  colorlinks   = true, %Colours links instead of ugly boxes
  urlcolor     = blue, %Colour for external hyperlinks
  linkcolor    = blue, %Colour of internal links
  citecolor   = red %Colour of citations
}

% Comments
\usepackage{comment}
\usepackage{todonotes}

% References
\usepackage{biblatex}
\addbibresource{references.bib}

%%%%%%%%%%%%%%%%%%%%%%%%%%%%%%%%%%%%%%%%%%%%%%%%%%%%%%%%%%%%%%%%%%%%%%%%%%%%%%%%

\title{Assignment 1 Deep Learning Group 33}

\author{Kian Niek Hamidi Abd Abad (Student ID), Márton Mohácsi (Student ID), Sid Pranjale (Student ID), Jorik Roebroek (Student ID)}

\date{ \today }

%%%%%%%%%%%%%%%%%%%%%%%%%%%%%%%%%%%%%%%%%%%%%%%%%%%%%%%%%%%%%%%%%%%%%%%%%%%%%%%%

\begin{document}
\maketitle

\tableofcontents

\noindent\fbox{
  \begin{minipage}{\dimexpr\textwidth-2\fboxsep-2\fboxrule\relax}
    \textbf{Assignment:} 
    \begin{enumerate}[(a)]
      \item Select your choice of neural networks model that is suitable for this task and motivate it. Train your model to predict one step ahead data point, during training (see following Figure). Scale your data before training and scale them back to be able to compare your predictions with real measurements.
      \item How many past time steps should you input into your network to achieve the best possible performance? (Hint: This is a tunable parameter and needs to be tuned).
      \item Once your model is trained, use it to predict the next 200 data points recursively. This means feeding each prediction back into the model to generate the subsequent predictions.
      \item On May 8th, download the real test dataset and evaluate your model by reporting both the Mean Absolute Error (MAE) and Mean Squared Error (MSE) between its predictions and the actual test values. Add
    \end{enumerate}
  \end{minipage}
}

\section{Choice of Neural Networks Model}
We have chosen to implement the [Insert Name of Neural Network Model] for our deep learning assignment. This model is well-suited for [briefly explain why this model is appropriate for the task at hand]. The architecture of this model allows it to effectively capture [mention specific features or patterns relevant to the task], making it a strong candidate for achieving high performance on our dataset.


\end{document}